\appendix
\setstretch{2} % triple spacing
\section{Proof of Section~\ref{sec:vanilla_mean}}

\subsection{Proof of Equation~\eqref{eq:ppi-mean-var}}
Since the unlabeled covariates $\{\widetilde{X}_j\}_{j=1}^N$ are i.i.d. from the same distribution $P$ as the labeled data, and $f$ is fixed with respect to the evaluation data, it follows that $\widetilde{f}_j = f(\widetilde{X}_j)$ are i.i.d. with
\[\operatorname{Var}\!\left(\widetilde{f}_j\right) = \operatorname{Var}\left(f(X)\right)=\sigma_f^2,
\quad
\mathbb{E}\!\left[\widetilde{f}_j\right] =\mathbb{E}[f(X)]= \mu_f.\]

\noindent Moreover, under the mild regularity condition that $\wt{f}_j$ has finite variance $\sigma_f^2$, we have\[\frac{\sqrt{N}\left(\bar f_U-\mu_f\right)}{\sigma_f}\overset{d}{\longrightarrow}\mathcal{N}(0,1),\quad\text{i.e.,}\quad \bar f_U\;\approx\;\mathcal{N}\left(\mu_f,\frac{\sigma_f^2}{N}\right)\;\text{for large $N$},\] by the Central Limit Theorem, where $\bar f_U=\frac{1}{N}\sum^N_{j=1}\wt{f}_j$. 

\shengyi{$\bar f_U\;=\;\mathcal{N}\left(\mu_f,\frac{\sigma_f^2}{N}\right)+o_p\!\left(\frac{1}{\sqrt{N}}\right)$}

\noindent In addition, since the labeled data $\{(X_i, Y_i)\}_{i=1}^n$ are i.i.d. from the distribution $P$, we have
\[\mathbb{E}\!\left[Y_i - f_i\right]
= \mathbb{E}[Y_i] - \mathbb{E}[f_i]
= \mathbb{E}[Y] - \mathbb{E}[f(X)]
= \theta^\star - \mu_f,
\quad \text{and} \quad
\operatorname{Var}\!\left(Y_i - f_i\right)
= \operatorname{Var}\!\left(Y - f(X)\right)
= \sigma_\varepsilon^2.\]

\noindent Similarly, under the mild regularity condition that $\sigma^2_\varepsilon$ is finite, we have
\[
\frac{\sqrt{n}\Bigg(\left(\overline{Y}_L- \bar f_L\right)-
\left(\theta^\star-\mu_f\right)\Bigg)}{\sigma_\varepsilon}\overset{d}{\longrightarrow}\mathcal{N}(0,1),
\quad\text{i.e.,}\quad 
\left(\overline{Y}_L- \bar f_L\right)\approx\mathcal{N}\left(\theta^\star-\mu_f,\frac{\sigma_\varepsilon^2}{n}\right)\;\text{for large $n$},
\] by the Central Limit Theorem, where $\overline{Y}_L =\frac{1}{n}\sum_{i=1}^n Y_i,\ \bar f_L=\frac{1}{n}\sum_{i=1}^n f_i$.

\shengyi{$\overline{Y}_L- \bar f_L\;=\;\mathcal{N}\left(\theta^\star-\mu_f,\frac{\sigma_\varepsilon^2}{n}\right)+o_p\!\left(\frac{1}{\sqrt{n}}\right)$}

\noindent Finally, under the independence of the labeled and unlabeled splits, and given~\eqref{eq:ppi-mean-est}, we have \[
\wh\theta_{\mathrm{PPI}}
\;=\;
\frac{1}{N}\sum_{j=1}^N \wt f_j
\;+\;\frac{1}{n}\sum_{i=1}^n \big(Y_i-f_i\big)=
\bar f_U+\left(\overline{Y}_L-\bar f_L\right)\approx\mathcal{N}\left(\theta^\star,\frac{\sigma_f^2}{N}+\frac{\sigma_\varepsilon^2}{n}\right)\;\text{for large $N$ and $n$},
\] as the sum of independent asymptotically normal variables is also asymptotically normal.

\shengyi{$\wh\theta_{\mathrm{PPI}}\;=\;\mathcal{N}\left(\theta^\star,\frac{\sigma_f^2}{N}+\frac{\sigma_\varepsilon^2}{n}\right)+o_p\!\left(\max\!\left\{
\frac{1}{\sqrt{N}},\,
\frac{1}{\sqrt{n}}
\right\}\right)$}
\subsection{Proof of Equation~\eqref{eq:n-for-ci-vanilla}}
\label{subsec:n-for-ci-vanilla}
\begin{equation*}
n \;\ge\; \frac{\sigma_\varepsilon^2}{(h/z_{1-\alpha/2})^2 - \sigma_f^2/N}
\quad\text{(feasible iff the denominator is positive).}
\end{equation*}

\yiqun{We have this conversation in slakc, but I think what you meant is for any N?}

\noindent Given the asymptotic result in~\eqref{eq:ppi-mean-var}, \[
\wh\theta_{\mathrm{PPI}} \approx \mathcal{N}\!\Bigg(\theta^\star,\ \frac{\sigma_f^2}{N} + \frac{\sigma_\varepsilon^2}{n}\Bigg)\quad\text{for large $N$ and $n$},
\]the asymptotic Wald $(1-\alpha)$ confidence interval for $\theta^\star$ is \[\widehat\theta_{\mathrm{PPI}}
\pm z_{1-\alpha/2}
\sqrt{\frac{\sigma_f^2}{N}+\frac{\sigma_\varepsilon^2}{n}}.\]

\noindent Therefore, for any fixed large $N$, the half-width of a $(1-\alpha)$ Wald confidence interval can be approximated as
\[
h\;\approx\;z_{1-\alpha/2}\sqrt{\frac{\sigma_f^2}{N} + \frac{\sigma_\varepsilon^2}{n}}.
\]
\yiqun{Why can you approximate that? I think we are skipping a very important step here}

\noindent
It follows that, to achieve a target half-width $h_t$, we require
\begin{align*}
z_{1-\alpha/2}\sqrt{\frac{\sigma_f^2}{N} + \frac{\sigma_\varepsilon^2}{n}} \;\le\; h_t
&\iff \frac{\sigma_f^2}{N} + \frac{\sigma_\varepsilon^2}{n} \;\le\; \left(\frac{h_t}{z_{1-\alpha/2}}\right)^2,\\[0.3em]
&\iff \frac{\sigma_\varepsilon^2}{n} \;\le\; \left(\frac{h_t}{z_{1-\alpha/2}}\right)^2 - \frac{\sigma_f^2}{N}.
\end{align*}

\noindent
In particular, if
\[
\left(\frac{h_t}{z_{1-\alpha/2}}\right)^2 - \frac{\sigma_f^2}{N} > 0,\]
then\[
n \;\ge\; \frac{\sigma_\varepsilon^2}{(h_t / z_{1-\alpha/2})^2 - \sigma_f^2 / N}.
\]
Otherwise, if
\[\left(\frac{h_t}{z_{1-\alpha/2}}\right)^2 - \frac{\sigma_f^2}{N} \le 0,\]
then $n < 0$, which is not feasible since $n$ denotes the labeled sample size and must be positive.

\noindent
Therefore, only when $\left(\tfrac{h_t}{z_{1-\alpha/2}}\right)^2 - \tfrac{\sigma_f^2}{N} > 0$, that is, when the predictions on the unlabeled data are not too noisy relative to the target precision, do we have a feasible range of $n$ in~\eqref{eq:n-for-ci-vanilla} that attains the target half-width of a $(1-\alpha)$ Wald confidence interval for $\theta^\star$.

\noindent
We summarize the above analysis in the following Lemma.

\begin{lemma}
\label{lem:n-CI}
Let $\hat{\theta}$ denote the estimator of the parameter of interest. If $\hat{\theta}$ is asymptotically normal, then to achieve a target half-width $h_t$ of a $(1-\alpha)$ Wald confidence interval corresponding to a two-sided test at level $\alpha$, the required sample size should satisfy
\[z_{1-\alpha/2}\sigma_{\hat \theta}\;\le\;h_t,\]
which is equivalently expressed as
\[\Var(\hat{\theta}) \;\le\; \left(\frac{h_t}{z_{1-\alpha/2}}\right)^2,\]subject to additional feasibility conditions when applicable.
\end{lemma}

\subsection{Proof of Equation~\eqref{eq:n-for-power-vanilla}}
\label{subsec:n-for-power-vanilla}
\begin{equation}
n \;\ge\; \frac{\sigma_\varepsilon^2}{\big(\Delta/(z_{1-\alpha/2}+z_{1-\beta})\big)^2 - \sigma_f^2/N}.
\tag{\ref{eq:n-for-power-vanilla}}
\end{equation}

\noindent
Let
\[
\sigma_{\mathrm{PPI}} = 
\sqrt{\frac{\sigma_f^2}{N} + \frac{\sigma_\varepsilon^2}{n}},
\]
denote the asymptotic standard deviation of $\widehat{\theta}_{\mathrm{PPI}}$ 
(as defined in~\eqref{eq:ppi-mean-var}).

\noindent We consider testing
\[
H_0:\ \theta^\star=\theta_0
\quad\text{versus}\quad
H_1:\ \theta^\star=\theta_0+\Delta,
\]
using the standardized test statistic
\[
Z_{stat}=\frac{\widehat{\theta}_{\mathrm{PPI}}-\theta_0}{
\sigma_{\mathrm{PPI}}},
\]where $\wh{\theta}_{\mathrm{PPI}}$ is an estimator of $\theta^\star$ satisfying the asymptotic normality
\[\wh{\theta}_{\mathrm{PPI}}
\approx
\mathcal{N}\!\Big(
\theta^\star,\,
\sigma_{\mathrm{PPI}}^2
\Big)
\quad\text{for large $N$ and $n$},\]as in~\eqref{eq:ppi-mean-var}.

\noindent
To achieve a desired power $1-\beta$ at significance level $\alpha$, we determine the required sample size $n$.

\yiqun{you should specify this test statistic in the proposition, and you should restate the $H_0$ here instead of later}

\noindent
Under the null hypothesis $ H_0\!:\theta^\star = \theta_0 $, by the asymptotic normality in~\eqref{eq:ppi-mean-var}, we have
\[
\widehat{\theta}_{\mathrm{PPI}} 
\approx 
\mathcal{N}\!\left(
  \theta_0,\ 
  \sigma_{\mathrm{PPI}}^2
\right),
\quad \text{i.e.,} \quad 
Z_{stat} \sim \mathcal{N}(0,1).
\]

\yiqun{Why is $\hat{\theta}_{PPI}$ approximately normal?}

\noindent
Therefore, for a two-sided test, the critical region is
\[
\left\{
\widehat{\theta}_{\mathrm{PPI}}:\ 
|Z_{stat}| \ge z_{1-\alpha/2}
\right\}.
\]

\noindent
Similarly, under the alternative hypothesis $ H_1\!:\theta^\star = \theta_0 + \Delta $, by the asymptotic normality in~\eqref{eq:ppi-mean-var}, we have
\[
\widehat{\theta}_{\mathrm{PPI}} 
\approx 
\mathcal{N}\!\left(
  \theta_0 + \Delta,\ 
  \sigma_{\mathrm{PPI}}^2
\right),
\quad \text{i.e.,} \quad 
Z_{stat} -
  \frac{\Delta}{\sigma_{\mathrm{PPI}}}
\sim 
\mathcal{N}(0,1).
\]

\yiqun{Why is $\hat{\theta}_{PPI} $approximately normal?}

\noindent
Hence, the power of the two-sided test is
\begin{align*}
    P_{H_1}\!\left(|Z_{stat}| \ge z_{1-\alpha/2}\right)&=
    P_{H_1}\!\left(Z_{stat} \ge z_{1-\alpha/2}\right) +
    P_{H_1}\!\left(Z_{stat} \le -z_{1-\alpha/2}\right)\\
    &=P_{H_1}\!\left(Z_{stat}-\frac{\Delta}{\sigma_{\mathrm{PPI}}} \ge z_{1-\alpha/2}-\frac{\Delta}{\sigma_{\mathrm{PPI}}}\right) +
    P_{H_1}\!\left(Z_{stat}-\frac{\Delta}{\sigma_{\mathrm{PPI}}} \le -z_{1-\alpha/2}-\frac{\Delta}{\sigma_{\mathrm{PPI}}}\right)\\
    &=1-\Phi\left(z_{1-\alpha/2}-\frac{\Delta}{\sigma_{\mathrm{PPI}}}\right)+\Phi\left(-z_{1-\alpha/2}-\frac{\Delta}{\sigma_{\mathrm{PPI}}}\right)\\
    &=\Phi\left(-z_{1-\alpha/2}+\Bigg|\frac{\Delta}{\sigma_{\mathrm{PPI}}}\Bigg|\right)+\Phi\left(-z_{1-\alpha/2}-\Bigg|\frac{\Delta}{\sigma_{\mathrm{PPI}}}\Bigg|\right)
\end{align*} where $\Phi(\cdot)$ denotes the standard normal cumulative distribution function.
\yiqun{For the derivations above, we write $P_{H_1}$ or $P_{H_0}$ since we are talking about frequentist probability here, we don't condition on the hypotheses}


\yiqun{why is the following claim true?}
\noindent Since
\[
\Phi\!\left(-z_{1-\alpha/2}
  + \Bigg|\frac{\Delta}{\sigma_{\mathrm{PPI}}}\Bigg|\right)
\ge
\Phi\!\left(-z_{1-\alpha/2}
  - \Bigg|\frac{\Delta}{\sigma_{\mathrm{PPI}}}\Bigg|\right),
\]
to achieve a power of $1 - \beta$, it suffices to require
\yiqun{are you saying that this is not an exact power analysis?}
\[
\Phi\!\left(-z_{1-\alpha/2}
  + \Bigg|\frac{\Delta}{\sigma_{\mathrm{PPI}}}\Bigg|\right)
\ge 1 - \beta.
\]

\noindent It follows that \[-z_{1-\alpha/2}+\Bigg|\frac{\Delta}{\sigma_{\mathrm{PPI}}}\Bigg|\;\ge\; z_{1-\beta},\quad\text{i.e.,}\quad z_{1-\alpha/2}+z_{1-\beta}\;\le\; \Bigg|\frac{\Delta}{\sigma_{\mathrm{PPI}}}\Bigg|.\]


\noindent There are two possible cases, 
\begin{enumerate}
  \item[(1)] $|z_{1-\alpha/2} + z_{1-\beta}|\;\le\; \Big|\frac{\Delta}{\sigma_{\mathrm{PPI}}}\Big|$
  \item[(2)] $z_{1-\alpha/2} + z_{1-\beta}\;\le\; -\Big|\frac{\Delta}{\sigma_{\mathrm{PPI}}}\Big|$
\end{enumerate}
\yiqun{I don' think you case (2) is correct. $a<|b|$ should imply $a<b; (b>0)$ or $a<-b \; (b<0)$? }

\noindent
In case~$(1)$, we have 
\begin{align*}
    \left(z_{1-\alpha/2} + z_{1-\beta}\right)^2 \;\le\; \frac{\Delta^2}{\sigma_{\mathrm{PPI}}^2}\;=\;
    \frac{\Delta^2}{\frac{\sigma_f^2}{N} + \frac{\sigma_\varepsilon^2}{n}}\;\le\;
    \frac{\Delta^2}{\frac{\sigma_f^2}{N}},
    \quad\text{i.e.,}\quad \frac{\Delta^2}{\left(z_{1-\alpha/2} + z_{1-\beta}\right)^2}-\frac{\sigma_f^2}{N}\;\ge\;0.
\end{align*}Moreover,
\begin{align*}
     \left(z_{1-\alpha/2} + z_{1-\beta}\right)^2 \;\le\; 
    \frac{\Delta^2}{\frac{\sigma_f^2}{N}
    + \frac{\sigma_\varepsilon^2}{n}}\;&\iff\;
    \frac{\sigma_\varepsilon^2}{n}\;\le\;\frac{\Delta^2}{\left(z_{1-\alpha/2} + z_{1-\beta}\right)^2}-\frac{\sigma_f^2}{N}\\
    &\iff\; n \;\ge\; \frac{\sigma_\varepsilon^2}{\Bigg(\Delta/(z_{1-\alpha/2}+z_{1-\beta})\Bigg)^2 - \sigma_f^2/N}
\end{align*}
In case~$(2)$, we have 
\begin{align*}
    z_{1-\alpha/2} + z_{1-\beta}
    \;\le\;
    -\Bigg|\frac{\Delta}{\sigma_{\mathrm{PPI}}}\Bigg|
    \;\le\; 0.
\end{align*}However, this violates the assumption of non-trivial hypothesis testing, as it contradicts the monotonicity of the inverse quantile function of standard normal distribution. Therefore, we only consider case~$(1)$.

\noindent
Therefore, to achieve the desired power $1-\beta$ at significance level $\alpha$, the labeled sample size $n$ must satisfy the condition in case~(1), as derived in~\eqref{eq:n-for-power-vanilla}.

\noindent
We summarize the power analysis in the following lemma.
\begin{lemma}
\label{lem:n-power}
Let $\hat{\theta}$ denote the estimator of the parameter of interest. For a two-sided test at level $\alpha$ and target power $1-\beta$, the required sample size must satisfy 
\[|z_{1-\alpha/2} + z_{1-\beta}| \;\le\; \Bigg|\frac{\Delta}{\sigma_{\hat \theta}}\Bigg|,\]
which is equivalently expressed as
\[
\Var(\hat \theta) \;\le\; \frac{\Delta^2}{\big(z_{1-\alpha/2} + z_{1-\beta}\big)^2}.
\]\end{lemma}

\yiqun{this section is important -- could you move it to the main text? open a new subsection talking about the estimation}
\subsection{Estimators of $\sigma^2_f$ and $\sigma^2_\varepsilon$ in Equations~\eqref{eq:n-for-ci-vanilla}--\eqref{eq:n-for-power-vanilla}}
\label{subsec:estimator_sigma_f}

\yiqun{why are we using unbiased? can we use MLE? Can you add the variance for these estimators? I think for the variance esitmation to be consistent, is the 2nd moment conditon enough}
Since population quantities under $P$,
\[
\sigma_f^2 = \Var(\wt{f}_j),
\quad\text{and}\quad
\sigma_\varepsilon^2 = \Var(Y_i - f_i),
\]
are unknown in practice, we estimate them using their unbiased sample counterparts.

\noindent
That is, \[
\hat\sigma^2_f=\frac{1}{N-1}\sum^N_{j=1}\left(\wt f_j-\bar{f}_U\right)^2,
\]where $\bar{f}_U=\frac{1}{N}\sum^N_{j=1}\wt f_j$, as $\wt f_j$ are i.i.d. samples from the same distribution of $f(X)$.

\noindent
Similarly, \[
\hat\sigma^2_\varepsilon=\frac{1}{n-1}\sum^n_{i=1}\left(\left(Y_i-f_i\right)-\left(\overline{Y}_L-\bar f_L\right)\right)^2,
\]where $\overline{Y}_L=\frac{1}{n}\sum^n_{i=1} Y_i$ and $\bar{f}_L=\frac{1}{n}\sum^n_{i=1} f_i$, as $Y_i-f_i$ are i.i.d. samples of $Y-f(X)$.

\noindent
In the analysis of confidence intervals and statistical power, we substitute $\widehat{\sigma}_f^2$ and $\widehat{\sigma}_\varepsilon^2$ into~\eqref{eq:n-for-ci-vanilla}--\eqref{eq:n-for-power-vanilla} when computing the required labeled sample size.

\section{Proof of Section~\ref{sec:PPI++_mean}}
\subsection{Proof of Equation~\eqref{eq:ppi++-mean-var}}
\begin{equation}
\Var(\wh\theta_\lambda)
=
\frac{\sigma_Y^2}{n}
+\lambda^2\Big(\frac{\sigma_f^2}{N}+\frac{\sigma_f^2}{n}\Big)
-\frac{2\lambda}{n}\Cov(Y,f).
\tag{\ref{eq:ppi++-mean-var}}
\end{equation}
Given that $\wh\theta_\lambda \;=\; \overline Y_L \;+\; \lambda\,\big(\bar f_U - \bar f_L\big)$, we can decompose its variance as
\begin{align*}
\Var(\wh\theta_\lambda)
&=\Var\left(\overline{Y}_L\right)
+\Var\left(\lambda\,\big(\bar f_U - \bar f_L\big)\right)
+2\cdot\Cov\left(\overline{Y}_L,\lambda\,\big(\bar f_U - \bar f_L\big)\right)\\
&=\Var\left(\overline{Y}_L\right)
+\lambda^2\cdot\Var\left(\bar f_U - \bar f_L\right)
+2\lambda\cdot \Cov\left(\overline{Y}_L,\bar f_U\right)
-2\lambda\cdot \Cov\left(\overline{Y}_L,\bar f_L\right)
\end{align*}

\noindent
We now consider each term individually. 

\noindent
\textbf{(i)} Since $Y_i$ are i.i.d. from $P_Y$, it follows that
\begin{align*}
    \Var\left(\overline{Y}_L\right)
    =\Var\left(\frac{1}{n}\sum^n_{i=1}Y_i\right)
    =\frac{1}{n^2}\sum^n_{i=1}\Var(Y_i)
    =\frac{1}{n}\Var(Y)
    =\frac{\sigma_Y^2}{n}.
\end{align*}

\noindent
\textbf{(ii)} Since the labeled data $\{(X_i, Y_i)\}_{i=1}^n$ and the unlabeled data $\{\wt {X}_j\}_{j=1}^N$ are independent, and the predictor $f$ is fixed with respect to the evaluation data, it follows that $f_i = f(X_i)$ and $\wt{f}_j = f(\widetilde{X}_j)$ are independent.

\noindent
Moreover, since $\{(X_i, Y_i)\}_{i=1}^n$ are i.i.d. with $\Var\left(f( X_i)\right)=\sigma_f^2$ and $\{\wt {X}_j\}_{j=1}^N$ are i.i.d. from the same distribution $P_X$ with $\Var\left(f(\wt X_j)\right)=\sigma_f^2$, we have \begin{align*}
    \Var\left(\bar f_U-\bar f_L\right)
    =\Var\left(\bar f_U\right)+\Var\left(\bar f_L\right)
    =\frac{\sigma^2_f}{N}+\frac{\sigma_f^2}{n}
\end{align*}

\noindent
\textbf{(iii)} The cross term 
$2\lambda\,\Cov\left(\overline{Y}_L,\bar f_U\right)$ vanishes 
because the labeled and unlabeled splits are independent, i.e.\
$\Cov\left(\overline{Y}_L,\bar f_U\right)=0$.

\noindent
\textbf{(iv)} Finally, for the remaining covariance term, we expand
\begin{align*}
    \Cov(\overline{Y}_L,\bar f_L)
    &=\frac{1}{n^2}\Cov\!\left(\sum_{i=1}^n Y_i,\,\sum_{j=1}^n f_j\right)\\
    &=\frac{1}{n^2}\sum_{i=1}^n\Cov(Y_i,f_i)
      +\frac{1}{n^2}\sum_{i\neq j}\Cov(Y_i,f_j).
\end{align*}
Since the labeled samples $\{(X_i, Y_i)\}_{i=1}^n$ are i.i.d., 
the pairs $\{(Y_i, f_i)\}_{i=1}^n$ are also i.i.d. for the fixed predictor $f$, 
implying that $\Cov(Y_i,f_i)=\Cov(Y,f)$. Moreover, because $Y_i$ and $f(X_j)$ are independent for $i \neq j$, 
the cross-covariance terms vanish. It then follows that\[
    \Cov(\overline{Y}_L,\bar f_L)
    = \frac{1}{n^2}\sum_{i=1}^n \Cov(Y_i,f_i)
    = \frac{1}{n}\Cov(Y,f).
\]

\noindent
Substituting all components back into the variance decomposition yields the desired result in~\eqref{eq:ppi++-mean-var}.

\subsection{Proof of Equation~\eqref{eq:lambda-star-mean} and~\eqref{eq:ppi++-mean-var-opt}}
Given the variance of $\hat{\theta}_\lambda$ in~\eqref{eq:ppi++-mean-var}, we define \[
g(\lambda)\coloneqq\Var\left(\hat{\theta}_\lambda\right)
=\frac{\sigma_Y^2}{n}
+\lambda^2\Bigg(\frac{\sigma_f^2}{N}+\frac{\sigma_f^2}{n}\Bigg)
-\frac{2\lambda}{n}\Cov(Y,f)\]
Taking the derivative of $g(\lambda)$ with respect to $\lambda$, we obtain
\begin{align*}
    g^\prime(\lambda)= 2\lambda\Bigg(\frac{\sigma_f^2}{N}+\frac{\sigma_f^2}{n}\Bigg)-\frac{2}{n}\Cov(Y,f)
\end{align*}
Setting $g^\prime(\lambda)=0$ gives \[\lambda^\star=\frac{\Cov(Y,f)}{n\left(\frac{\sigma_f^2}{N}+\frac{\sigma_f^2}{n}\right)}=\frac{\Cov(Y,f)}{\sigma_f^2\left(1+\frac{n}{N}\right)}=\frac{\Cov(Y,f)}{\sigma_f^2(1+r)}\]where $r=\frac{n}{N}$. Since $g(\lambda)$ is quadratic in $\lambda$, it is convex and therefore attains its unique global minimum at $\lambda^\star$, which is the \emph{oracle} choice in~\eqref{eq:lambda-star-mean}.

\noindent
Substituting $\lambda^\star$ into~\eqref{eq:ppi++-mean-var}, we have\begin{align*}
    \Var\left(\hat{\theta}_{\lambda^\star}\right)
    &=\frac{\sigma^2_Y}{n}
    +\left(\frac{\Cov(Y,f)}{\sigma_f^2\left(1+\frac{n}{N}\right)}\right)^2\cdot\sigma_f^2\cdot\left(\frac{1}{N}+\frac{1}{n}\right)
    -\frac{2}{n}\cdot\frac{\Cov(Y,f)}{\sigma_f^2\left(1+\frac{n}{N}\right)}\cdot\Cov(Y,f)\\
    &=\frac{\sigma^2_Y}{n}+\frac{\Cov(Y,f)^2}{\sigma_f^2}\cdot\frac{N}{n(N+n)}-2\cdot \frac{\Cov(Y,f)^2}{\sigma_f^2}\cdot\frac{N}{n(N+n)}\\
    &=\frac{\sigma^2_Y}{n}-\frac{\Cov(Y,f)^2}{\sigma_f^2}\cdot\frac{N}{n(N+n)}
\end{align*}
Therefore, the minimized variance is \[\Var\left(\hat{\theta}_{\lambda^\star}\right)=\frac{\sigma^2_Y}{n}-\frac{\Cov(Y,f)^2}{\sigma_f^2}\cdot\frac{N}{n(N+n)}\]as in~\eqref{eq:ppi++-mean-var-opt}.

\subsection{Proof of Equation~\eqref{eq:n-for-ci-ppi++}}
Since $\overline{Y}_L$, $\bar f_U$, and $\bar f_L$ are sample means of random variables sharing the same mean and variance, they are asymptotically normal under the regularity condition that the variances are finite, i.e., $\Var(Y)=\sigma_Y^2<\infty$ and $\Var(f(X))=\Var(f(\wt X))=\sigma_f^2<\infty$, for large $N$ and $n$.

\noindent
Therefore, given the definition of \emph{PPI++} family in~\eqref{eq:ppi++-mean-est}, $\wh{\theta}_{\lambda^\star}$ is asymptotically normal for large $n$ and $N$ under standard regularity conditions, as it is a sum of asymptotically normal random variables.

\noindent
We first show that $\wt{\theta}_{\lambda^\star}$ is unbiased:\[\E\left(\wh{\theta}_{\lambda^\star}\right)
=\E\left[\overline Y_L \;+\; \lambda^\star\,\big(\bar f_U - \bar f_L\big)\right]
=\E\left[\overline Y_L\right] \;+\; \lambda^\star\left(\E\left[\,\bar f_U\right]-\; \E\left[\,\bar f_L\right]\right)
=\theta^\star+\lambda^\star(\mu_f-\mu_f)=\theta^\star.
\]

\noindent
Given $\Var\left(\hat{\theta}_{\lambda^\star}\right)$ in~\eqref{eq:ppi++-mean-var-opt}, we have \[
\hat{\theta}_{\lambda^\star}\;\approx\;\mathcal{N}\left(\theta^\star,\ \frac{\sigma^2_Y}{n}-\frac{\Cov(Y,f)^2}{\sigma_f^2}\cdot\frac{N}{n(N+n)}\right)\quad\text{for large $N$ and $n$}.
\]

\noindent
We then consider testing
\[
H_0:\ \theta^\star=\theta_0
\quad\text{versus}\quad
H_1:\ \theta^\star=\theta_0+\Delta,
\]
using the standardized test statistic
\[
Z=\frac{\wh{\theta}_{\lambda^\star}-\theta_0}{
\sigma_{\lambda^\star}},\]where $\sigma_{\lambda^\star}$ denotes the standard deviation of $\wh{\theta}_{\lambda^\star}$.

\noindent
Using Lemmas~\ref{lem:n-CI} and~\ref{lem:n-power}, the labeled sample size $n$ required to achieve either a target half-width of a $(1-\alpha)$ Wald confidence interval or a desired power $1-\beta$ can be obtained by solving
\[
\Var\!\left(\widehat{\theta}_{\lambda^\star}\right) \;\le\; S^2
\]for some specified threshold $S^2$.

\noindent
Substituting $\Var\left(\hat{\theta}_{\lambda^\star}\right)$ in~\eqref{eq:ppi++-mean-var-opt}, we have
\begin{align*}
    \frac{\sigma_Y^2}{n}-\frac{\Cov(Y,f)^2}{\sigma_f^2}\cdot\frac{N}{n(n+N)} \;\le\; S^2
    \iff S^2n^2+\left(S^2N-\sigma^2_Y\right)n+\left(\frac{\Cov(Y,f)^2}{\sigma_f^2}N-N\sigma^2_Y\right)\;\geq\;0
\end{align*}
\noindent
Let
\[
a = S^2,\quad
b = S^2 N - \sigma_Y^2,\quad
c = \frac{\Cov(Y,f)^2}{\sigma_f^2}N - N\sigma_Y^2.
\]
We can then rewrite the inequality as
\[a n^2 + b n + c \ge 0\]
\noindent
Since $a = S^2 > 0$, the inequality holds for $n\ge n_1$ or $n\le n_2$, where \[
n_{1,2} =\frac{-b\pm\sqrt{b^2-4ac}}{2a},
\]by the quadratic formula.

\noindent
Because the labeled sample size $n$ must be positive, the feasible region is $n \ge n_1$, provided that $n_1 > 0$ and $b^2 - 4ac \ge 0$.

\noindent
That is, \begin{equation*}
n \;\ge\; 
\frac{ \sigma_Y^2 - S^2 N + \sqrt{(\sigma_Y^2 - S^2 N)^2 + 4 S^2\left(\sigma_Y^2 N - \frac{\Cov(Y,f)^2}{\sigma_f^2}\,N\right)} }{2 S^2},
\end{equation*} as in~\eqref{eq:n-for-ci-ppi++}.

\noindent
And this result holds when both\begin{align*}
    \frac{ \sigma_Y^2 - S^2 N + \sqrt{(\sigma_Y^2 - S^2 N)^2 + 4 S^2\left(\sigma_Y^2 N - \frac{\Cov(Y,f)^2}{\sigma_f^2}\,N\right)} }{2 S^2}\;&>\;0,\\
    (\sigma_Y^2 - S^2 N)^2 + 4 S^2\left(\sigma_Y^2 N - \frac{\Cov(Y,f)^2}{\sigma_f^2}\,N\right)\;&\geq\;0.
\end{align*}

\subsection{Estimators of $\sigma_Y^2$ and $\Cov(Y,f)$ in Equation~\eqref{eq:ppi++-mean-var}}

The quantities $\sigma_Y^2$, $\sigma_f^2$, and $\Cov(Y,f)$ are unknown population parameters appearing in the variance of $\widehat{\theta}_\lambda$, as well as in its minimized variance and the subsequent analyses of confidence intervals and statistical power.

\noindent
We have already introduced the unbiased estimator of $\sigma_f^2$ in Section~\ref{subsec:estimator_sigma_f}.
In this section, we focus on the estimation of $\sigma_Y^2$ and $\Cov(Y,f)$.

\noindent
In practice, we replace $\sigma_Y^2$ with its unbiased sample estimator,\[
\wh{\sigma}_Y^2
= \frac{1}{n - 1}
\sum_{i = 1}^n
\left(Y_i - \overline{Y}_L\right)^2,\]
where
$\overline{Y}_L = \frac{1}{n}\sum_{i = 1}^n Y_i$, since the $\{Y_i\}_{i=1}^n$ are i.i.d. draws from the population distribution of $P_Y$.

\noindent
Next, since the pairs $\{\left(Y_i,f_i\right)\}^n_{i=1}$ are i.i.d. for the fixed predictor $f$, we estimate $\Cov(Y,f)$ using the sample covariance,\[
\wh \Cov(Y,f) = \frac{1}{n-1}\sum^n_{i=1}\left(Y_i-\overline{Y}_L\right)\left(f_i-\bar f_L\right),
\]where $\overline{Y}_L=\frac{1}{n}\sum^n_{i=1} Y_i$ and $\bar{f}_L=\frac{1}{n}\sum^n_{i=1} f_i$.

\section{Proof of Section~\ref{sec:t-test}}
\subsection{Proof of Equation~\eqref{eq:ppi++-delta-mean-var}}

Given the definition of $\wh\Delta_{\mathrm{PPI}}$ in~\eqref{eq:ppi++-dealta-mean}, we have
\[\wh\Delta_{\mathrm{PPI}}={\wh{\theta}_{A,\mathrm{PPI}}}-{\wh{\theta}_{B,\mathrm{PPI}}},\]where ${\wh{\theta}_{A,\mathrm{PPI}}}=\overline{\wt f}_A - \overline{ f}_A + \overline{Y}_A$, and ${\wh{\theta}_{B,\mathrm{PPI}}}=\overline{\wt f}_B - \overline{ f}_B + \overline{Y}_B$.

\noindent
Analogously to the \emph{vanilla PPI} mean estimator for a single group in~\eqref{eq:ppi-mean-est}, under independence of labeled and unlabeled splits within each group and under mild regularity,
\begin{equation*}
\wh\theta_{A,\mathrm{PPI}} \approx \mathcal{N}\!\left(\theta_A^\star,\ \frac{\sigma_{f,A}^2}{N_A} + \frac{\sigma_{\varepsilon,A}^2}{n_A}\right),\quad
\wh\theta_{B,\mathrm{PPI}} \approx \mathcal{N}\!\left(\theta_B^\star,\ \frac{\sigma_{f,B}^2}{N_B} + \frac{\sigma_{\varepsilon,B}^2}{n_B}\right),
\end{equation*}for large $N_A$, $n_A$, $N_B$, and $n_B$.

\noindent
Furthermore, since groups $A$ and $B$ are independent, it follows that\[
\wh\theta_{A,\mathrm{PPI}}-\wh\theta_{B,\mathrm{PPI}} \approx \mathcal{N}\!\left(\theta_A^\star-\theta_B^\star,\ \frac{\sigma_{f,A}^2}{N_A} + \frac{\sigma_{\varepsilon,A}^2}{n_A}+\ \frac{\sigma_{f,B}^2}{N_B} + \frac{\sigma_{\varepsilon,B}^2}{n_B}\right)
\]and hence,\[
\Var\left(\wh\theta_{A,\mathrm{PPI}}-\wh\theta_{B,\mathrm{PPI}}\right)\approx
\frac{\sigma_{f,A}^2}{N_A} + \frac{\sigma_{\varepsilon,A}^2}{n_A}+\ \frac{\sigma_{f,B}^2}{N_B} + \frac{\sigma_{\varepsilon,B}^2}{n_B},
\]for large $N_A$, $n_A$, $N_B$, and $n_B$.

\section{Proof of Section~\ref{sec:cross-fitting}}
\subsection{Proof of Equation~\eqref{eq:cross-ppi++-distrn}}
The proof follows a similar template as the proof of Theorem~$2$ in~\citep{zrnic2024cross} and combines the two lemmas used therein.

\begin{lemma}\label{lem:asy-fU} (Lem.~3 in the proof of Thm.~2 in~\citep{zrnic2024cross})

\noindent
Suppose that the predictions are stable (Ass.~\ref{ass:M-stable-f}). Denote
\[\tilde{L}_j= \frac{1}{\sqrt{N}} \sum_{i=1}^N\Big(
\nabla \tilde{\ell}^{f^{(j)}}_{\theta,i}
- \mathbb{E}[\nabla \ell^{f^{(j)}}_\theta]
- (\nabla \tilde{\ell}_{\theta,i}^{\bar f}- \mathbb{E}[\nabla \ell_\theta^{\bar f}])\Big).\]
Then,
\[\frac{1}{K} \sum_{j=1}^K \tilde{L}_j \;\xrightarrow{p}\;0.\]
\end{lemma}

\begin{lemma}\label{lem:asy-fL} (Lem.~4 in the proof of Thm.~2 in~\citep{zrnic2024cross})

\noindent
Suppose that the predictions are stable (Ass.~\ref{ass:M-stable-f}). Denote
\[L_j= \frac{1}{\sqrt{n}} \sum_{i \in I_j}
\Big(\nabla \ell^{f^{(j)}}_{\theta,i}
- \mathbb{E}[\nabla \ell^{f^{(j)}}_\theta]
- (\nabla \ell_{\theta,i}^{\bar f}- \mathbb{E}[\nabla \ell_\theta^{\bar f}])\Big).\]
Then,
\[\sum_{j=1}^K L_j \;\xrightarrow{p}\; 0.\]
\end{lemma}

\noindent
Writing \begin{align*}\nabla L(\theta)&=\E[\ell_\theta]\\
    &=\frac{1}{K}\left(\sum^K_{j=1}\E\left[\lambda \ell_\theta^{f^{(j)}}\right]-\E\left[\lambda \ell_\theta^{f^{(j)}}-l_\theta\right]\right)\\
    &=\frac{1}{KN}\sum^K_{j=1}\sum^N_{i=1}\E\left[\lambda \ell_\theta^{f^{(j)}}\right]
    -\frac{1}{n}\sum^K_{j=1}\sum_{i\in I_j}\E\left[\lambda \ell_\theta^{f^{(j)}}-\ell_\theta\right]
    \end{align*}where $n$ is assumed to be divisible by $K$ for simplicity, that is, $|I_j|=\frac{n}{K}$ for all $j$.

\noindent
We have\begin{align*}
&\sqrt{n}\,
\Big( \tfrac{n}{N}\,\lambda^2\,\bar{\Sigma}_\theta + \bar{\Sigma}_{\Delta_\lambda,\theta} \Big)^{-1/2}
\big( \nabla L_\lambda^+(\theta) - \nabla L_\lambda(\theta) \big) \\
=&
\sqrt{n}\,
\Big( \tfrac{n}{N}\,\lambda^2\,\bar{\Sigma}_\theta + \bar{\Sigma}_{\Delta_\lambda,\theta} \Big)^{-1/2}
\left(
\frac{1}{KN} \sum_{j=1}^K \sum_{i=1}^N
\nabla \lambda\tilde{\ell}^{f^{(j)}}_{\theta,i}
\;-\;
\frac{1}{n} \sum_{j=1}^K \sum_{i \in I_j}
\big( \nabla \lambda\ell^{f^{(j)}}_{\theta,i} - \nabla \ell_{\theta,i} \big)
\;-\;
\nabla L(\theta)
\right) \\
=&
\sqrt{n}\,
\Big( \tfrac{n}{N}\,\lambda^2\,\bar{\Sigma}_\theta + \bar{\Sigma}_{\Delta_\lambda,\theta} \Big)^{-1/2}\\
&
\left(
\frac{1}{KN} \sum_{j=1}^K \sum_{i=1}^N
\big(
\nabla \lambda \tilde{\ell}^{f^{(j)}}_{\theta,i}
-
\mathbb{E}[\nabla \lambda \ell^{f^{(j)}}_\theta]
\big)
\;-\;
\frac{1}{n} \sum_{j=1}^K \sum_{i \in I_j}
\big(
(\nabla \lambda \ell^{f^{(j)}}_{\theta,i} - \nabla \ell_{\theta,i})
-
\mathbb{E}[\nabla \lambda \ell^{f^{(j)}}_\theta - \nabla \ell_\theta]
\big)
\right) \\
=&
\sqrt{n}\,
\Big( \tfrac{n}{N}\,\lambda^2\,\bar{\Sigma}_\theta + \bar{\Sigma}_{\Delta_\lambda,\theta} \Big)^{-1/2}
\big( T_{1,\lambda} - T_{2,\lambda} \big),
\end{align*}where we define \begin{align*}
T_{1,\lambda}&=\frac{1}{KN} \sum_{j=1}^K \sum_{i=1}^N
\big(\nabla \lambda \tilde{\ell}^{f^{(j)}}_{\theta,i}
-\mathbb{E}[\nabla  \lambda \ell^{f^{(j)}}_\theta]\big),\\
T_{2,\lambda}&=\frac{1}{n} \sum_{j=1}^K \sum_{i \in I_j}
\big((\nabla  \lambda \ell^{f^{(j)}}_{\theta,i} - \nabla \ell_{\theta,i})
-\mathbb{E}[\nabla  \lambda \ell^{f^{(j)}}_\theta - \nabla \ell_\theta]\big).
\end{align*}

\noindent
In particular, we can write
\begin{align*}
\sqrt{N} T_{1,\lambda}
&= \frac{1}{\sqrt{N}} \sum_{i=1}^N
\big( \nabla \lambda \tilde{\ell}_{\theta,i}^{\bar f} - \mathbb{E}[\nabla \lambda \ell^{\bar f}_\theta] \big)
+\frac{\lambda}{K} \sum_{j=1}^K\tilde{L}_j\\
& =\frac{1}{\sqrt{N}} \sum_{i=1}^N
\big( \nabla \lambda \tilde{\ell}_{\theta,i}^{\bar f}
- \mathbb{E}[\nabla \lambda \ell^{\bar f}_\theta] \big)
+ o_P(1),\\
\sqrt{n} T_{2,\lambda} &= 
  \frac{1}{\sqrt{n}} \sum_{j=1}^K \sum_{i \in I_j}\Big( \nabla \lambda \ell^{f^{(j)}}_{\theta,i} - \nabla \ell_{\theta,i}
  - \mathbb{E}[\nabla \lambda \ell^{\bar f}_\theta - \nabla \ell_\theta] \Big)
  + \lambda \sum_{j=1}^K L_j\\
  &=\frac{1}{\sqrt{n}} \sum_{j=1}^K \sum_{i \in I_j}\Big( \nabla \lambda \ell^{f^{(j)}}_{\theta,i} - \nabla \ell_{\theta,i}
  - \mathbb{E}[\nabla \lambda \ell^{\bar f}_\theta - \nabla \ell_\theta] \Big)
  +o_P(1)
\end{align*}using Lemma~\ref{lem:asy-fU} and Lemma~\ref{lem:asy-fL}, respectively.

\noindent
Therefore, \begin{align*}
    &\sqrt{n}\,
\Big( \tfrac{n}{N}\,\lambda^2\,\bar{\Sigma}_\theta + \bar{\Sigma}_{\Delta_\lambda,\theta} \Big)^{-1/2}
\big( \nabla L_\lambda^+(\theta) - \nabla L_\lambda(\theta) \big)\\ 
=&
\sqrt{n}\,
\Big( \tfrac{n}{N}\,\lambda^2\,\bar{\Sigma}_\theta + \bar{\Sigma}_{\Delta_\lambda,\theta} \Big)^{-1/2}
\big( T_{1,\lambda} - T_{2,\lambda} \big)\\
=&\Big( \tfrac{n}{N}\,\lambda^2\,\bar{\Sigma}_\theta + \bar{\Sigma}_{\Delta_\lambda,\theta} \Big)^{-1/2}
\left(\frac{\sqrt{n}}{\sqrt{N}}\ \sqrt{N}\  T_{1,\lambda} - \sqrt{n}T_{2,\lambda} \right)\\
=&\Big( \tfrac{n}{N}\,\lambda^2\,\bar{\Sigma}_\theta + \bar{\Sigma}_{\Delta_\lambda,\theta} \Big)^{-1/2}\frac{\sqrt{n}}{\sqrt{N}}\frac{1}{\sqrt{N}} \sum_{i=1}^N
\big( \nabla \lambda \tilde{\ell}_{\theta,i}^{\bar f}
- \mathbb{E}[\nabla \lambda \ell^{\bar f}_\theta] \big)\\
&-\Big( \tfrac{n}{N}\,\lambda^2\,\bar{\Sigma}_\theta + \bar{\Sigma}_{\Delta_\lambda,\theta} \Big)^{-1/2}\frac{1}{\sqrt{n}} \sum_{j=1}^K \sum_{i \in I_j}\Big( \nabla \lambda \ell^{f^{(j)}}_{\theta,i} - \nabla \ell_{\theta,i}
  - \mathbb{E}[\nabla \lambda \ell^{\bar f}_\theta - \nabla \ell_\theta] \Big)
\end{align*}

Denote the limits of $\frac{n}{N}$, $\bar\Sigma_\theta$, $\bar\Sigma_{\Delta_\lambda,\theta}$ by $r$, $\bar\Sigma_{\theta,\lim}$, $\bar\Sigma_{\Delta_\lambda,\theta,\lim}$, respectively.

By the Lindeberg central limit theorem, the first term converge in distribution to\[
    \mathcal{N}\left(0,
    \Big( r\,\lambda^2\,\bar{\Sigma}_{\theta,\lim} + \bar{\Sigma}_{\Delta_\lambda,\theta,\lim} \Big)^{-1/2}
    \left(r\lambda^2\bar{\Sigma}_{\theta,\lim}\right)
    \Big( r\,\lambda^2\,\bar{\Sigma}_{\theta,\lim} + \bar{\Sigma}_{\Delta_\lambda,\theta,\lim} \Big)^{-1/2}
    \right),
\]and the second term converge in distribution to\[
    \mathcal{N}\left(0,
    \Big( r\,\lambda^2\,\bar{\Sigma}_{\theta,\lim} + \bar{\Sigma}_{\Delta_\lambda,\theta,\lim} \Big)^{-1/2}
    \bar{\Sigma}_{\Delta_\lambda,\theta,\lim} 
    \Big( r\,\lambda^2\,\bar{\Sigma}_{\theta,\lim} + \bar{\Sigma}_{\Delta_\lambda,\theta,\lim} \Big)^{-1/2}
    \right),
\]respectively, as $\Big( \tfrac{n}{N}\,\lambda^2\,\bar{\Sigma}_\theta + \bar{\Sigma}_{\Delta_\lambda,\theta} \Big)^{-1/2}$ is symmetric.

\noindent
Since the labeled and unlabeled samples are independently drawn, the two terms are independent. Therefore, we can sum their expectations and variances, obtaining\begin{align*}
    \sqrt{n}\,
\Big( \tfrac{n}{N}\,\lambda^2\,\bar{\Sigma}_\theta + \bar{\Sigma}_{\Delta_\lambda,\theta} \Big)^{-1/2}
\big( \nabla L_\lambda^+(\theta) - \nabla L_\lambda(\theta) \big)\overset{d}{\longrightarrow}\mathcal{N}(0,I)
\end{align*}as in~\eqref{eq:cross-ppi++-distrn}.


